\documentclass{article}

\usepackage[a4paper, left=1.5cm, right=1.5cm, top=2cm, bottom=2cm]{geometry}

\usepackage{../../../components/mathtex}

\usepackage{hyperref}
\usepackage{fancyhdr}
\usepackage{ccicons}
\usepackage{caption}
\usepackage{subcaption}
\usepackage{algorithm}
\usepackage{algorithmic}

% Configuration des en-têtes et pieds de page
\pagestyle{fancy}
\fancyhf{} % reset tout

\fancyhead[L]{UFR MATHS INFO}
\fancyhead[C]{Assemblée générale}
\fancyhead[R]{2026-2027}

\fancyfoot[L]{Ewen Rodrigues de Oliveira\\\ccbyncsa}
\fancyfoot[R]{\thepage}

\begin{document}

\docTitle{AG UFR Math-Info 6 novembre}

Merci à Igor et Loïc pour leur aide dans la prise de notes.
C'est un résumé prise de parole par prise de parole.

\disclaimer{De mon point de vue, \textbf{rien n'a été décidé à cette AG}.\\
Tout le monde est d'accord pour dire qu'il faut une confiance.\\
L'UFR demande un retour en cours pour pouvoir donner des notes, et \textbf{semble} opposé au "10 améliorable".\\\\

Pour un résumé complet, passez ça dans Gemini ou ChatGPT vous savez faire. Désolé pour les fautes d'orthographe, le résumé n'a pas encore été relu.\\
Un planning hebdomadaire va être mis en place pour savoir quels bâtiments seront ouverts et quand.}

\section{Ordre du jour}

\begin{enumerate}
    \item Point sur la semaine passée
    \item Résolution des problèmes lié au service santé étudiant
    \item Revendication étudiants
    \item MCC
    \item Informations sur la suite du blocus
\end{enumerate}

\section{Point sur la semaine passée}
\begin{itemize}
    \item Incompréhension entre profs et présents à l'AG, c'était bien un débrayage. Erreur de communication.\\ 
    Conseil du doyen de la faculté des sens : écrire les revendications avant le vote.
    \item Personne qui a été huée parce que anti-blocus. Tout le monde devrait pouvoir parler librement, et appel à ceux qui sont contre les blocus à parler "plus fort" et que "leur avis".
    \item La tribune de l'agée a rappelé à l'ordre les personnes qui ont huées et les personnes opposées au blocus ont pu continuer à parler. La personne invite tout le monde à participer.
    \item La personne trouve que le vote n'est pas secret et que la personne doit pouvoir voter sans être jugée. Elle veut davantage de discussion et de représentation. ELle veut savoir à quoi sert cette AG.
    \item L'objectif de l'AG est d'avoir un contact entre les enseignants et les élèves. Ce n'est pas à cette AG de discuter du blocage. Par contre, discuter sur les MCC et l'avis de chacun sur Sophie Germain. Il pense que l'AG n'est pas inutile. Idée d'élire un représentant pour porter la voix de l'UFR à l'AG principale.
    \item Il y a des personnes mobilisées qui ne bloquent pas. L'AG intra-UFR sert à parler aux enseignants. Certains enseignant disent qu'il n'y a pas de communication entre les enseignants et les élèves. L'AG est un lieu de discussion.
    \item Les écrans 
    \item inquietude des etudiants sur les MCC et les examens, notamment la penalisation de l'absence des CC avec les blocus.
    \item Beaucoup de detresse morale et psychologiques, peu de communications eleves profs. Mettre ressources sur Moodle, certains MI L1 veulent des ressources.
    \item Les L3 veulent aider les L1 et L2, mais il y a un manque de communication. Des ressources sont sur ewenrdo.fr et essayer de prendre un lien.
    \item Mathrix (@mathrixup) pour le tutorat, en présentiel à Bibliothèque François Mitterrand.
    \item Décider d'un représentant pour voir ce qui sera communiqué à l'AG principale.
    \item Il y a plusieurs moyens de bloquer. On emmerde pas le gouvernement, mais les étudiants pour sûr. Risque de décrochage, problèmes pour ceux qui ne sont pas supportés par leurs parents. Ouverture de la fac pour réunir les étudiants, former, etc. Pas normal que certains profs aient mis des examens à la BU.
    \item Ceux qui sont mobilisés pensent pas que les personnes qui bloquent sont des horribles personnes. Tous ne sont pas d'accord pour le blocus.

\end{itemize}

\section{Pôle santé}

Problème : \textbf{Le SSE est dans Sophie Germain qui est bloqué.}
Inacceptable. Le SSE est ouvert, il faut passer par la porte de derrière.

\section{Revendications étudiants}
\begin{itemize}
    \item les bloqueurs rappellent qu'ils luttent pour nos droits, face a la decrepitation de la situation économique des universités, le manque flagrant de moyens
    \item ils appellent a la solidarité des UFR, qui envoie des mails, qui participent aux Universités populaires
    \item une sage de 24 ans nous rappellent que la vie n'est pas finie avec les blocus
    \item \textbf{il n'y a pas de pénalisation pour les absences, même les boursiers. Mais ils ne peuvent pas l'écrire.}
    \item L'ouverture partielle permet d'arrêter le distanciel.
    \item C'est faux. À Condorcet ils ont fait les deux.
    \item On est pas dans une démocratie fonctionnelle. Le gouvernement n'en a pas rien à faire. Forte répression chez les lycéens. La grève est un bon moyen.
    \item Il n'y a pas une "minorité de bloqueurs". Ceux qui ont raison c'est ceux qui ont froid en hiver, etc. Ceux qui viennent pas parce qu'ils sont loin peuvent quand même être solidaires.
    \item Peut être proposer le choix en bis : visio et présentiel en même temps. Sans contrôle d'assiduité.
    \item Il faudrait ouvrir les salles SANS CC.
    \item Doyen : les conseils d'UFR soutiennent totalement les revendications, même sur la répression par la violence à des revendications légitimes. Il y a une distance entre les étudiants et les enseignants. La confiance doit revenir. Il faut revenir à la fac, même si c'est pas pour faire cours. Si un prof fait un CC, le doyen prend l'engagement que ça ne sera pas pris en compte.
    \item demande une intervention de responsables d'UFR pour eviter que certains controles est lieu
    \item communication entre les UFR, pour le relais handicap
    \item les HaF pas adaptés pour les étudiants photosensibles comme ce neuil de Loïc 
\end{itemize}

\section{MCC}    
\begin{itemize}
    \item on veut des cours en ligne, les cc en ligne a faire pour tester chez soit (pas noté), communiquer aux étudiants les modalités et les attendus des prochains CC, demande de limitation des contenus present aux examens finales, demander plus d'annales sur moodle, plus facile de faire le tutorat a la BU, 
    \item demande de faire les cours en ligne l'apres midi.
    \item proposition de notes planchers. il ne faut pas que le distanciel soit dans les examens.
    \item "note amélioration" : même note qu'on pourrait améliorer pour ceux qui le veulent.
    \item cours en distanciel : injuste. isolement social.
    \item faire des DM. - pire idée 
    \item Non possibilité de faire de cours en présentiel si aucune organisation.
    \item debrayage semaine prochaine a Halle aux farines, avec un edt qui permettrait de savoir quand et comment la fac serait ouverte.
    \item apparament note plancher impossible par le tribunal administratif (take du prof d'info qui tient pas trop la route)
    \item CC sur les 4 premieres semaines.
    \item idée : math en université populaire.
    \item jury doit rester souverain. pas de généralité pour le 10 améliorable
    \item grève des surveillants d'examen.
    \item L3 doivent avoir une note ! il faut revenir en cours pour avoir des notes. reprendre les cours. 
    \item les bonnes conditions pour avoir des notes, il n'y a pas de conditions pour avoir des notes.
\end{itemize}

\textbf{Doyen :} Les notes sont nécessaires pour obtenir le diplôme. Mais les programmes seront modifiées. Dans une même promotion, les étudiants doivent avoir les mêmes conditions, dans chaque cursus. Les nouvelles MCC devront être données 15 jours avant les examens, et doivent être à l'avantage des élèves.

\begin{itemize}
    \item les enseignants devraient rejoindre concrètement la lutte. Le "10 +" est pas légal, mais c'est pas avec des mesures légales qu'on va obtenir quelque chose.
    \item Il n'y a pas que nous qui allons avoir ces problématiques. Idem pour les lycéens par exemple et le bac.
    \item il faudrait avoir, même en master, des compensations entre cours, en particulier dans le même bloc.
    \item les notes ne sont pas un reflet de ce qu'on a compris, c'est faux. Le système de note est pas bon.
    \item les ufr ne répondent pas à la chose, juste répondre "il faut retourner" en cours. 
    \item les profs sont pas méchants, ils veulent la réussite des élèves.
    \item voir filière par filière pour les notes. il faut une notation adaptée dans chaque cours.
    \item inquiétude pour ceux qui veulent faire un master de tout ce qui est perdu, et qui ne sera pas rattrapable.
    \item les étudiants veulent des engagements des enseignants, mais pour l'instant il n'y en a pas.
    \item les étudiants ont besoin d'une sécurité pour les examens, surtout pour les profs qui ne sont pas solidaires.
    \item Vassout : les étudiants ont majoritairement demandé du distanciel. Les profs restent mécontents de ce système. "on est pas là pour vous mettre des batons dans les roues". il faut revenir en présenciel pour avoir des cours et matière à évaluer.
    \item les notes sont un tri social, et empêchent d'accéder à un meilleur avenir. 
\end{itemize}

\section{Informations sur la suite du blocus}

\begin{itemize}
    \item Un planning hebdomadaire va être mis en place pour savoir quels bâtiments sont bloqués et quand.
\end{itemize}

\end{document}
